\documentclass[11pt,letterpaper]{article}

% --------------------------------------------------
% Packages
% --------------------------------------------------
\usepackage{amsmath,amssymb,amsthm,mathtools}
\usepackage{geometry}
\usepackage{setspace}
\usepackage{csquotes}
\usepackage{hyperref}
\usepackage{bm}

% --------------------------------------------------
% Page setup
% --------------------------------------------------
\geometry{margin=1in}
\setstretch{1.15}

% --------------------------------------------------
% Theorem environments
% --------------------------------------------------
\newtheorem{definition}{Definition}
\newtheorem{theorem}{Theorem}
\newtheorem{proposition}{Proposition}
\newtheorem{corollary}{Corollary}

% --------------------------------------------------
% Title information
% --------------------------------------------------
\title{Intelligence as Compression:\\
Toward a Theory of Mind Recognition via Effective Assembly}
\author{Fkyxion}
\date{\today}

\begin{document}

\maketitle

\begin{abstract}
This essay proposes that intelligence is best understood not as problem-solving capacity, symbolic manipulation, or behavioral adaptability per se, but as \emph{compression under constraint}. Building on assembly-based measures of complexity and integrating recent work on semantic infrastructure, operational mereology, and field-based theories of cognition, I introduce the \emph{Effective Assembly Index for Mind Recognition} (EAIMR). The EAIMR formalizes mind-likeness as a ratio between irreducible constructive depth and the structural affordances that permit reuse, parallelism, degeneracy, and coherence. The resulting framework offers a principled criterion for recognizing minds across biological, artificial, and hybrid substrates without appealing to anthropocentric or behavioral tests.
\end{abstract}

\section{Introduction}

The question \enquote{What is intelligence?} has traditionally been answered by appeal to outward behavior, internal representations, or normative rationality. Each of these approaches presupposes a privileged ontology: agents, symbols, goals, or environments. Yet across biology, computation, and physics, intelligence appears wherever systems \emph{compress structure while remaining dynamically responsive}.

This essay advances the thesis that intelligence is not a substance or faculty but a \emph{compression regime}. Minds are not defined by what they know or what they do, but by how efficiently they assemble, reuse, and stabilize structure across scales. Intelligence, on this view, is a thermodynamically and semantically constrained process of structural condensation.

To make this precise, I introduce a quantitative proposal for mind recognition grounded in assembly theory and semantic coherence: the Effective Assembly Index for Mind Recognition (EAIMR).

\section{Compression Beyond Information Theory}

Shannon information theory formalizes compression as entropy reduction relative to a code. However, Shannon entropy is indifferent to meaning, construction history, and reuse. Two strings with identical entropy may differ radically in the effort required to produce them or the semantic structure they support.

Compression relevant to intelligence must therefore be:
\begin{enumerate}
  \item \textbf{Constructive}: sensitive to how structure is assembled.
  \item \textbf{Reusable}: capable of supporting repeated function.
  \item \textbf{Composable}: able to integrate parts without collapse.
  \item \textbf{Coherent}: maintaining global consistency under transformation.
\end{enumerate}

These criteria motivate a shift from purely statistical compression to \emph{assembly-based} measures.

\section{Assembly as a Measure of Cognitive Depth}

The \emph{assembly index} $A(x)$ of an object or process $x$ measures the minimal number of distinct operations required to construct $x$ from basic primitives, allowing reuse of intermediate components. Assembly differs from Kolmogorov complexity in that it is operational, historical, and sensitive to material and semantic constraints.

In cognitive systems, assembly captures:
\begin{itemize}
  \item developmental depth,
  \item architectural differentiation,
  \item learned modular structure.
\end{itemize}

However, raw assembly alone overcounts brute accumulation. A pile of unrelated components may have high assembly but little intelligence. Intelligence requires that assembly be \emph{effectively organized}.

\section{Effective Assembly and Mind Recognition}

To distinguish mere complexity from mind-like organization, we define the Effective Assembly Index for Mind Recognition.

\begin{definition}[EAIMR]
For a system or process $x$, the Effective Assembly Index for Mind Recognition is defined as
\[
\mathrm{EAIMR}(x) =
\frac{A(x)}{R(x)\,P(x)\,D(x)} \cdot C(x),
\]
where:
\begin{itemize}
  \item $A(x)$ is the assembly index of $x$,
  \item $R(x)$ measures reuse,
  \item $P(x)$ measures parallelism,
  \item $D(x)$ measures degeneracy,
  \item $C(x)$ measures coherence.
\end{itemize}
\end{definition}

Each denominator term discounts assembly by structural efficiencies, while coherence reweights the result toward integrated systems.

\subsection{Reuse}

Reuse $R(x)$ quantifies the degree to which components or substructures are invoked multiple times across the system. High reuse indicates compression via abstraction: the same structure performs multiple roles without duplication.

\subsection{Parallelism}

Parallelism $P(x)$ measures the extent to which assemblies can occur or operate simultaneously. Parallel systems compress time and energy by exploiting independence without fragmentation.

\subsection{Degeneracy}

Degeneracy $D(x)$ captures the availability of distinct structures capable of fulfilling similar functions. Unlike redundancy, degeneracy supports robustness and adaptive compression by allowing functional equivalence across heterogeneous components.

\subsection{Coherence}

Coherence $C(x)$ measures the degree to which the system maintains global consistency across transformations, contexts, or perturbations. Coherence distinguishes integrated minds from disjoint toolkits.

\section{Intelligence as a Compression Regime}

Under EAIMR, intelligence corresponds to a specific region in structural space:
\begin{quote}
High assembly that cannot be trivially reduced, discounted by maximal reuse, parallelism, and degeneracy, yet amplified by coherence.
\end{quote}

Brute-force systems inflate $A(x)$ without corresponding gains in $R$, $P$, or $D$, yielding low EAIMR. Simple automata exhibit high reuse but low assembly. Minds occupy the narrow corridor where all terms are simultaneously nontrivial.

\begin{proposition}
A system exhibits mind-like organization if and only if $\mathrm{EAIMR}(x)$ exceeds a context-relative threshold while remaining stable under perturbation.
\end{proposition}

This reframes mind recognition as a structural inference problem rather than a behavioral test.

\section{Relation to Field-Based and Semantic Theories of Mind}

Within field-based theories such as scalar--vector--entropy frameworks, assembly corresponds to structured condensation of field configurations, reuse to attractor reuse, parallelism to distributed flow, degeneracy to basin overlap, and coherence to topological invariants.

Operational mereology further grounds EAIMR by replacing static part-whole relations with event-sourced assembly histories. Minds are not sets of components but replayable construction logs stabilized by compression.

\section{Implications}

The EAIMR framework has several consequences:
\begin{enumerate}
  \item \textbf{Substrate neutrality}: minds may exist in biological, artificial, or hybrid systems.
  \item \textbf{Gradualism}: mind-likeness is scalar, not binary.
  \item \textbf{Anti-anthropocentrism}: intelligence does not require language, goals, or self-models.
  \item \textbf{Detectability}: mind recognition becomes an empirical question about structure and assembly.
\end{enumerate}

\section{Conclusion}

Intelligence is compression that survives contact with reality. The Effective Assembly Index for Mind Recognition formalizes this intuition by tying mind-likeness to the irreducible depth of construction discounted by structural efficiencies and amplified by coherence. Rather than asking whether a system behaves like a mind, we may instead ask whether it \emph{assembles the world efficiently enough to matter}.

This shift reframes intelligence as a property of structure, not appearance—and opens the possibility of recognizing minds wherever compression learns to hold itself together.

\end{document}

